%----------------------------------------------------------------------------------------
% SOMFORMULES. Alle opgaven zijn eigen werk.
%
% Opbouw: Start, leerstof (kern, uitgewerkt voorbeeld, korte oefening), en daarna op een
% nieuwe pagina de Opdrachten (door elkaar, stijgende moeilijkheid), Koppel en Test jezelf.
% Labels van dit hoofdstuk beginnen met h3-.
%----------------------------------------------------------------------------------------

\mychapter{3}{Somformules}

\lkr{Leerplan B+S'': LPD 44, formule voor de som van de eerste $n$ termen, met samengestelde intrest als context. Twee bewijzen die leerlingen moeten kunnen volgen en navertellen: de som tweemaal schrijven, en $q \cdot s_n - s_n$. De rode draad in de opdrachten: eerst herkennen welke soort rij het is, dan het aantal termen bepalen, dan pas de formule gebruiken. Oneindige sommen komen niet in dit hoofdstuk: die vragen de limiet uit hoofdstuk 4.}

\begin{oefeningen}[Start]
\begin{oefening}
\blokkop{Opwarmer}
\begin{enumerate}
  \item Een rekenkundige rij heeft $u_1 = 4$ en $v = 3$. Bereken $u_{10}$.
  \item Een meetkundige rij heeft $u_1 = 2$ en $q = 3$. Bereken $u_5$.
  \item Hoeveel termen telt de rij $8,\ 11,\ 14,\ \ldots,\ 77$?
  \item Los op: $n^2 + 3n - 40 = 0$.
  \item Werk uit: $(x-1)(x^3 + x^2 + x + 1)$.
  \item Los op: $2^x = 1024$ en $3^x = 200$.
\end{enumerate}
\begin{solution}
\textbf{a.}~31 \quad \textbf{b.}~162 \quad \textbf{c.}~24 \quad \textbf{d.}~$n=5$ of $n=-8$ \quad \textbf{e.}~$x^4-1$ \quad \textbf{f.}~$x=10$; $x\approx 4{,}82$
\end{solution}
\begin{fullsolution}
\textbf{a.}~$u_{10}=4+9\cdot 3=31$.\\
\textbf{b.}~$u_5=2\cdot 3^4=162$.\\
\textbf{c.}~$8+(n-1)\cdot 3=77$ geeft $n-1=23$, dus $n=24$.\\
\textbf{d.}~$(n-5)(n+8)=0$.\\
\textbf{e.}~$x^4+x^3+x^2+x-x^3-x^2-x-1=x^4-1$: bijna alles valt weg.\\
\textbf{f.}~$2^{10}=1024$, dus $x=10$. $x=\log_3 200=\dfrac{\log 200}{\log 3}\approx 4{,}82$.
\end{fullsolution}
\end{oefening}

\begin{oefening}\label{oef:h3-verken}
\blokkop{Verken}\quad Je spaart tien weken lang. In week 1 leg je \eur 3 opzij, en elke week \eur 2 meer dan de week ervoor.
\begin{enumerate}
  \item Schrijf de tien bedragen op een rij. Schrijf daaronder dezelfde bedragen in omgekeerde volgorde.
  \item Tel per kolom de twee getallen op. Wat valt op?
  \item Hoeveel heb je na tien weken gespaard? Gebruik de twee rijen, zonder alle bedragen op te tellen.
\end{enumerate}
\begin{solution}
\textbf{b.}~elke kolom geeft 24 \quad \textbf{c.}~\eur 120
\end{solution}
\begin{fullsolution}
\textbf{a.}~$3,\ 5,\ 7,\ 9,\ 11,\ 13,\ 15,\ 17,\ 19,\ 21$ en daaronder $21,\ 19,\ \ldots,\ 3$.\\
\textbf{b.}~Bovenaan komt er telkens 2 bij, onderaan gaat er telkens 2 af: elke kolom geeft 24.\\
\textbf{c.}~De twee rijen samen zijn $10\cdot 24=240$. Dat is twee keer het gespaarde bedrag, dus \eur 120.
\end{fullsolution}
\end{oefening}
\end{oefeningen}

\lkr{Laat oefening \ref{oef:h3-verken} maken vóór de uitleg. Vraag waarom je deelt door 2: dat is de stap die leerlingen het vaakst vergeten als ze de formule later moeten verklaren.}

%========================================================================================
\section{Som van een rekenkundige rij}

\subsection{De somformule}

\begin{theorie}
\begin{minipage}[t]{.64\linewidth}
$s_n$ is de som van de eerste $n$ termen:

\qquad $\formule{s_n = u_1 + u_2 + \cdots + u_n}$

\vspace{1.5mm}
Bij een rekenkundige rij is dat $n$ keer het gemiddelde van de eerste en de laatste term:

\vspace{1mm}
\qquad $\formule{s_n = n \cdot \dfrac{u_1 + u_n}{2}}$

\vspace{1.5mm}
De figuur toont het idee: twee keer dezelfde rij vormt een rechthoek van $n$ op $u_1 + u_n$.
\end{minipage}\hfill
\begin{minipage}[t]{.32\linewidth}
\vspace{-1mm}
\begin{tikzpicture}[x=5.4mm,y=2.7mm,font=\scriptsize]
  \foreach \i/\h in {0/2,1/3,2/4,3/5,4/6}{
    \fill[accent!45] (\i,0) rectangle (\i+1,\h);
    \draw[accent] (\i,0) rectangle (\i+1,\h);
    \draw[accent] (\i,\h) rectangle (\i+1,8);}
  \draw[decorate,decoration={brace,mirror,amplitude=3pt},grijs] (0,-.3) -- (5,-.3) node[midway,below=2pt,text=black] {$n$ termen};
  \draw[decorate,decoration={brace,mirror,amplitude=3pt},grijs] (5.15,0) -- (5.15,8) node[midway,right=2pt,text=black] {$u_1 + u_n$};
\end{tikzpicture}
\end{minipage}

\begin{bewijs}
\begin{stappen}
s_n = u_1 + u_2 + \cdots + u_n & De som. \\
s_n = u_n + u_{n-1} + \cdots + u_1 & Dezelfde som, omgekeerd. \\
2 s_n = (u_1 + u_n) + \cdots + (u_n + u_1) & Optellen per kolom. Bovenaan komt er telkens $v$ bij, onderaan gaat er telkens $v$ af: elke kolom geeft $u_1 + u_n$. \\
2 s_n = n \,(u_1 + u_n) & Er zijn $n$ kolommen. \\[1ex]
s_n = n \cdot \dfrac{u_1 + u_n}{2} & Delen door 2.
\end{stappen}
\end{bewijs}
\end{theorie}

\voorbeeld{}{Bereken $s_{20}$ van de rij $7,\ 11,\ 15,\ \ldots$\label{vb:h3-rr}}{%
\begin{stappen}
v = 11 - 7 = 4 & Rekenkundig. Controle: $15 - 11 = 4$. \\
u_{20} = 7 + 19 \cdot 4 = 83 & Eerst de laatste term: 19 sprongen. \\[1ex]
s_{20} = 20 \cdot \dfrac{7 + 83}{2} = 900 & Twintig keer het gemiddelde, en dat is 45.
\end{stappen}}

\nujij{Bereken $s_{30}$ van de rij $20,\ 18,\ 16,\ \ldots$}{$v = -2$, $u_{30} = -38$, $s_{30} = 30 \cdot \tfrac{20 - 38}{2} = -270$.}

\thuis{Vul je $u_n = u_1 + (n-1)v$ in, dan krijg je een tweede vorm: $s_n = \tfrac{n}{2}\bigl(2u_1 + (n-1)v\bigr)$. Die is handig als je $u_n$ nog niet kent.}

\begin{oefeningen}
\begin{oefening}
Bereken $s_{25}$ van de rij $9,\ 13,\ 17,\ \ldots$
\begin{solution}
1425
\end{solution}
\begin{fullsolution}
$v=4$, dus $u_{25}=9+24\cdot 4=105$ en $s_{25}=25\cdot\dfrac{9+105}{2}=1425$.
\end{fullsolution}
\end{oefening}
\end{oefeningen}

\subsection{Eerst het aantal termen}

\begin{werkwijze}
Is de laatste term gegeven en niet het aantal termen, bereken dan eerst $n$ met het expliciet voorschrift. Het aantal termen is één meer dan het aantal sprongen.
\end{werkwijze}

\voorbeeld{}{Bereken de som van alle veelvouden van 6 tussen 100 en 400.\label{vb:h3-aantal}}{%
\begin{stappen}
102 + 108 + \cdots + 396 & Eerst uitschrijven: $102 = 6 \cdot 17$ en $396 = 6 \cdot 66$. \\
102 + (n-1) \cdot 6 = 396 \ \Rightarrow \ n = 50 & 49 sprongen, dus 50 termen. \\[1ex]
s_{50} = 50 \cdot \dfrac{102 + 396}{2} = 12\,450 &
\end{stappen}}

\nujij{Bereken $50 + 54 + 58 + \cdots + 250$.}{$n = 51$, en de som is $51 \cdot 150 = 7650$.}

\begin{oefeningen}
\begin{oefening}
Bereken $15 + 22 + 29 + \cdots + 204$.
\begin{solution}
3066
\end{solution}
\begin{fullsolution}
$15+(n-1)\cdot 7=204$ geeft $n-1=27$, dus $n=28$. Dan is de som $28\cdot\dfrac{15+204}{2}=3066$.
\end{fullsolution}
\end{oefening}
\end{oefeningen}

\subsection{Het aantal termen uit de som}

\begin{werkwijze}
Is de som gegeven, schrijf dan $s_n$ met $n$ als onbekende. Dat geeft een vierkantsvergelijking. Enkel een oplossing in $\mathbb{N}_0$ is een aantal termen.
\end{werkwijze}

\voorbeeld{}{Hoeveel termen van $3,\ 7,\ 11,\ \ldots$ moet je optellen om 210 te krijgen?\label{vb:h3-nuit}}{%
\begin{stappen}
u_n = 3 + (n-1) \cdot 4 = 4n - 1 & De laatste term, met $n$ als onbekende. \\[1ex]
s_n = n \cdot \dfrac{3 + 4n - 1}{2} = n\,(2n + 1) & Invullen in de somformule. \\
2n^2 + n - 210 = 0 & $s_n = 210$ stellen. \\[1ex]
n = 10 \quad\text{of}\quad n = -\dfrac{21}{2} & Enkel $n = 10$ kan. Controle: $10 \cdot 21 = 210$.
\end{stappen}}

\nujij{Hoeveel termen van $6,\ 8,\ 10,\ \ldots$ moet je optellen om 150 te krijgen?}{$s_n = n(n+5)$, dus $n = 10$.}

\begin{oefeningen}
\begin{oefening}
Hoeveel termen van $4,\ 7,\ 10,\ \ldots$ moet je optellen om 175 te krijgen?
\begin{solution}
10
\end{solution}
\begin{fullsolution}
$u_n=3n+1$, dus $s_n=n\cdot\dfrac{4+3n+1}{2}=\dfrac{n(3n+5)}{2}$. Uit $3n^2+5n-350=0$ volgt $n=10$ of $n=-\tfrac{35}{3}$. Enkel $n=10$ kan.
\end{fullsolution}
\end{oefening}
\end{oefeningen}

%========================================================================================
\section{Som van een meetkundige rij}

\begin{oefeningen}
\begin{oefening}\label{oef:h3-mverken}
\blokkop{Verken}\quad Bij sneeuw verwittigt de school 1 leerling. Die belt 3 andere leerlingen. Elk van hen belt er weer 3, enzovoort. Noem $S = 1 + 3 + 9 + 27 + 81$ het aantal verwittigde leerlingen na vier belrondes.
\begin{enumerate}
  \item Schrijf ook $3S$ uit, als som van vijf termen.
  \item Bereken $3S - S$. Wat valt weg?
  \item Bereken daarmee $S$ zonder de termen op te tellen.
\end{enumerate}
\begin{solution}
\textbf{a.}~$3+9+27+81+243$ \quad \textbf{b.}~$243-1=242$ \quad \textbf{c.}~121
\end{solution}
\begin{fullsolution}
\textbf{a.}~Elke term maal 3: $3S=3+9+27+81+243$.\\
\textbf{b.}~De termen 3, 9, 27 en 81 staan in beide sommen en vallen weg: $3S-S=243-1=242$.\\
\textbf{c.}~$2S=242$, dus $S=121$.
\end{fullsolution}
\end{oefening}
\end{oefeningen}

\lkr{Laat de leerlingen vergelijken met de spaaroefening: bij een rekenkundige rij kies je voor optellen in omgekeerde volgorde, bij een meetkundige rij voor vermenigvuldigen met $q$. Vraag waarom de truc van de ene niet werkt bij de andere.}

\subsection{De somformule}

\begin{theorie}
Voor een meetkundige rij met $q \neq 1$: \qquad $\formule{s_n = u_1 \cdot \dfrac{q^{\,n} - 1}{q - 1}}$

\vspace{1.5mm}
Voor $q = 1$ zijn alle termen gelijk: \quad $\formule{s_n = n \cdot u_1}$

\begin{bewijs}
\begin{stappen}
s_n = u_1 + u_1 q + \cdots + u_1 q^{\,n-1} & De som. \\
q\, s_n = u_1 q + \cdots + u_1 q^{\,n-1} + u_1 q^{\,n} & Vermenigvuldig met $q$: alles schuift een plaats op. \\
q\, s_n - s_n = u_1 q^{\,n} - u_1 & Aftrekken: de termen ertussen vallen weg. \\
s_n (q - 1) = u_1 (q^{\,n} - 1) & Links en rechts buiten haakjes brengen. \\[1ex]
s_n = u_1 \cdot \dfrac{q^{\,n} - 1}{q - 1} & Delen door $q - 1$, wat mag omdat $q \neq 1$.
\end{stappen}
\end{bewijs}
\end{theorie}

Let op: in de somformule staat de exponent $n$, het aantal termen. In het expliciet voorschrift staat $n - 1$, het aantal sprongen.

\voorbeeld{}{Bereken $s_{10}$ van de rij $3,\ 6,\ 12,\ \ldots$\label{vb:h3-mr}}{%
\begin{stappen}
q = \dfrac{6}{3} = 2 & Meetkundig. Controle: $\dfrac{12}{6} = 2$. \\[1.5ex]
s_{10} = 3 \cdot \dfrac{2^{10} - 1}{2 - 1} = 3069 & Exponent 10, want er zijn 10 termen.
\end{stappen}}

\nujij{Bereken $s_9$ van de rij $1,\ {-2},\ 4,\ \ldots$}{$q = -2$, dus $s_9 = \tfrac{(-2)^9 - 1}{-2 - 1} = \tfrac{-513}{-3} = 171$.}

\begin{oefeningen}
\begin{oefening}
Bereken $s_6$ van de rij $2,\ 10,\ 50,\ \ldots$
\begin{solution}
7812
\end{solution}
\begin{fullsolution}
$q=5$, dus $s_6=2\cdot\dfrac{5^6-1}{5-1}=2\cdot\dfrac{15\,624}{4}=7812$.
\end{fullsolution}
\end{oefening}
\end{oefeningen}

\subsection{Eerst het aantal termen}

\begin{werkwijze}
Is de laatste term gegeven, bereken dan eerst $n$ uit $u_n = u_1 \cdot q^{\,n-1}$, zo nodig met een logaritme. Werk met de exacte machten, niet met afgeronde getallen.
\end{werkwijze}

\voorbeeld{}{Bereken $2 + 6 + 18 + \cdots + 4374$.\label{vb:h3-mrn}}{%
\begin{stappen}
2 \cdot 3^{\,n-1} = 4374 \ \Rightarrow \ 3^{\,n-1} = 2187 & $q = 3$. Delen door 2. \\
n - 1 = 7 \ \Rightarrow \ n = 8 & Want $3^7 = 2187$. Dus 8 termen. \\[1ex]
s_8 = 2 \cdot \dfrac{3^8 - 1}{3 - 1} = 6560 &
\end{stappen}}

\nujij{Bereken $7 + 14 + 28 + \cdots + 448$.}{$2^{n-1} = 64$, dus $n = 7$, en de som is $7 \cdot (2^7 - 1) = 889$.}

\begin{oefeningen}
\begin{oefening}
Bereken $5 + 15 + 45 + \cdots + 3645$.
\begin{solution}
5465
\end{solution}
\begin{fullsolution}
$5\cdot 3^{n-1}=3645$ geeft $3^{n-1}=729=3^6$, dus $n=7$. Dan is $s_7=5\cdot\dfrac{3^7-1}{3-1}=5465$.
\end{fullsolution}
\end{oefening}
\end{oefeningen}

%========================================================================================
\section{Het somteken}

\subsection{Lezen en berekenen}

\begin{theorie}
\begin{minipage}[t]{.56\linewidth}
Een som schrijf je kort met het \emph{somteken} $\Sigma$:

\vspace{1mm}
\qquad $\formule{\displaystyle\sum_{i=1}^{n} u_i = u_1 + u_2 + \cdots + u_n}$

\vspace{1.5mm}
De letter $i$ is de \emph{lopende index}: die neemt elk geheel getal aan van de beginwaarde tot en met de eindwaarde.
\end{minipage}\hfill
\begin{minipage}[t]{.4\linewidth}
\textbf{Aantal termen:} eindwaarde $-$ beginwaarde $+ 1$

\vspace{1.5mm}
\textbf{Rekenregels:}

$\formule{\displaystyle\sum (u_i + w_i) = \sum u_i + \sum w_i}$

$\formule{\displaystyle\sum c \cdot u_i = c \cdot \sum u_i}$
\end{minipage}
\end{theorie}

\thuis{Je leest $\sum_{i=1}^{n} u_i$ als: de som van de $u_i$ voor $i$ van 1 tot $n$. Het teken is de Griekse hoofdletter sigma, de S van \emph{som}. De naam van de index maakt niet uit: $\sum_{i=1}^{5} i$ en $\sum_{k=1}^{5} k$ zijn hetzelfde getal.}

\voorbeeld{}{Bereken $\displaystyle\sum_{i=3}^{12} (2i + 1)$ en $\displaystyle\sum_{j=1}^{5} 3 \cdot 2^{\,j}$.\label{vb:h3-somteken}}{%
\begin{stappen}
7 + 9 + 11 + \cdots + 25 & Vul $i = 3$, 4, 5 en 12 in: rekenkundig. \\
12 - 3 + 1 = 10 \text{ termen} & Niet 12 en niet 9. \\[1ex]
10 \cdot \dfrac{7 + 25}{2} = 160 & De eerste som. \\
6 + 12 + 24 + 48 + 96 & Vul $j = 1$, 2, 3 in: meetkundig, met $u_1 = 6$ en niet 3. \\[1ex]
6 \cdot \dfrac{2^5 - 1}{2 - 1} = 186 & De tweede som.
\end{stappen}}

\lkr{Laat leerlingen altijd de eerste drie termen uitschrijven voor ze een formule kiezen. Dat voorkomt de twee klassieke fouten: het verkeerde aantal termen, en $u_1$ aflezen zonder invullen.}

\nujij{Bereken $\displaystyle\sum_{k=2}^{9} (5k - 3)$.}{Termen $7,\ 12,\ \ldots,\ 42$; 8 termen; de som is $8 \cdot \tfrac{7 + 42}{2} = 196$.}

\begin{oefeningen}
\begin{oefening}
Bereken.
\begin{reeks}[2]
  \item $\displaystyle\sum_{i=2}^{11} (3i + 1)$
  \item $\displaystyle\sum_{j=1}^{6} 2 \cdot 3^{\,j}$
\end{reeks}
\begin{solution}
\textbf{a.}~205 \quad \textbf{b.}~2184
\end{solution}
\begin{fullsolution}
\textbf{a.}~$7+10+\cdots+34$: 10 termen, dus $10\cdot\dfrac{7+34}{2}=205$.\\
\textbf{b.}~$6+18+\cdots$: $u_1=6$, $q=3$ en 6 termen, dus $6\cdot\dfrac{3^6-1}{3-1}=2184$.
\end{fullsolution}
\end{oefening}
\end{oefeningen}

\subsection{De som als rij}

\begin{theorie}
De sommen $s_1,\ s_2,\ s_3,\ \ldots$ vormen zelf een rij. Daaruit vind je de termen terug:

\vspace{1mm}
\qquad $\formule{u_1 = s_1} \quad\text{en}\quad \formule{u_n = s_n - s_{n-1}} \quad\text{voor } n \geqslant 2$
\end{theorie}

\voorbeeld{}{Gegeven: $s_n = n^2 + 2n$. Bepaal $u_n$. Welke soort rij is dat?\label{vb:h3-uitsom}}{%
\begin{stappen}
u_1 = s_1 = 3 & Vervang $n$ door 1. \\
s_{n-1} = (n-1)^2 + 2(n-1) = n^2 - 1 & Vervang $n$ door $n - 1$. \\
u_n = n^2 + 2n - (n^2 - 1) = 2n + 1 & Voor $n \geqslant 2$. Ook $u_1 = 3$ past in dit voorschrift. \\
\text{rekenkundig met } v = 2 & Controle: $3 + 5 + 7 = 15 = s_3$.
\end{stappen}}

\begin{oefeningen}
\begin{oefening}
Gegeven: $s_n = n^2 + 4n$. Bepaal $u_n$.
\begin{solution}
$u_n=2n+3$
\end{solution}
\begin{fullsolution}
$s_{n-1}=(n-1)^2+4(n-1)=n^2+2n-3$, dus $u_n=s_n-s_{n-1}=2n+3$ voor $n\geqslant 2$. Ook $u_1=s_1=5$ past hierin.
\end{fullsolution}
\end{oefening}
\end{oefeningen}

%========================================================================================
\section{Sparen met vaste stortingen}

Je stort aan het einde van elk jaar een bedrag $a$ op een rekening met $p\%$ samengestelde intrest per jaar. Elke storting groeit apart aan, met $q = 1 + \tfrac{p}{100}$. De laatste storting brengt nog niets op.

\vspace{2mm}
\begin{tikzpicture}[x=2.25cm,font=\small]
  \draw[->,lijngrijs] (.55,0) -- (5.6,0) node[right,text=black] {tijd};
  \foreach \x/\t/\w in {1/{jaar 1}/{a\,q^{\,n-1}},2/{jaar 2}/{a\,q^{\,n-2}},4/{jaar $n-1$}/{a\,q},5/{jaar $n$}/{a}}{
    \fill[accent] (\x,0) circle (2pt);
    \node[above=1pt,font=\scriptsize,text=grijs] at (\x,0) {\t};
    \node[below=2pt] at (\x,0) {$\w$};}
  \node[fill=white,inner sep=1pt] at (3,0) {$\cdots$};
\end{tikzpicture}

\begin{theorie}
De eindwaarde na $n$ jaar is de som van een meetkundige rij met eerste term $a$: \quad $\formule{E = a \cdot \dfrac{q^{\,n} - 1}{q - 1}}$
\end{theorie}

\lkr{Dit deel past het best bij Economie-Wiskunde. Het leerplan noemt enkel samengestelde intrest als context: sla de opdrachten over sparen over als je dit deel niet behandelt. Leningen komen alleen voor in een opdracht op niveau 3.}

\voorbeeld{}{Je stort 6 jaar lang aan het einde van elk jaar \eur 1000 op een rekening met 3\% intrest. Hoeveel staat er na de laatste storting op? Hoeveel daarvan is intrest?\label{vb:h3-sparen}}{%
\begin{stappen}
q = 1 + \dfrac{3}{100} = 1{,}03 & \\[1.5ex]
E = 1000 \cdot \dfrac{1{,}03^6 - 1}{1{,}03 - 1} \approx 6468{,}41 & Na 6 jaar: ongeveer \textbf{\eur 6468,41}. \\[1.5ex]
6468{,}41 - 6 \cdot 1000 = 468{,}41 & \textbf{\eur 468,41} is intrest.
\end{stappen}}

\nujij{Je stort 10 jaar lang aan het einde van elk jaar \eur 500 op een rekening met 2\% intrest. Hoeveel staat er na de laatste storting op?}{$500 \cdot \tfrac{1{,}02^{10} - 1}{0{,}02} \approx 5474{,}86$ euro.}

\begin{oefeningen}
\begin{oefening}
Je stort 8 jaar lang aan het einde van elk jaar \eur 1500 op een rekening met 2\% intrest. Hoeveel staat er na de laatste storting op?
\begin{solution}
Ongeveer \eur 12\,874,45
\end{solution}
\begin{fullsolution}
$E=1500\cdot\dfrac{1{,}02^8-1}{0{,}02}\approx 12\,874{,}45$.
\end{fullsolution}
\end{oefening}
\end{oefeningen}

%----------------------------------------------------------------------------------------
% OPDRACHTEN: door elkaar. Niveau 1, dan 2, dan 3.
%----------------------------------------------------------------------------------------
\opdrachtenkop

\lkr{De opdrachten staan bewust door elkaar, met stijgende moeilijkheid. In deze versie staat bij elke opdracht het onderwerp, zodat je kunt kiezen.}

\begin{oefeningen}
\begin{oefening}[1]\label{oef:h3-soort}
\thema{gemengd: soort rij herkennen}Is de rij rekenkundig of meetkundig? Bereken dan de som.
\begin{enumerate}
  \item $4 + 9 + 14 + \cdots$ (12 termen)
  \item $3 + 6 + 12 + \cdots$ (7 termen)
  \item $20 + 17 + 14 + \cdots$ (10 termen)
  \item $2 - 6 + 18 - \cdots$ (6 termen)
\end{enumerate}
\begin{solution}
\textbf{a.}~rekenkundig, 378 \quad \textbf{b.}~meetkundig, 381 \quad \textbf{c.}~rekenkundig, 65 \quad \textbf{d.}~meetkundig, $-364$
\end{solution}
\begin{fullsolution}
\textbf{a.}~$v=5$, $u_{12}=59$, dus $12\cdot\dfrac{4+59}{2}=378$.\\
\textbf{b.}~$q=2$, dus $3\cdot(2^7-1)=381$.\\
\textbf{c.}~$v=-3$, $u_{10}=-7$, dus $10\cdot\dfrac{20-7}{2}=65$.\\
\textbf{d.}~$q=-3$, dus $2\cdot\dfrac{(-3)^6-1}{-3-1}=2\cdot\dfrac{728}{-4}=-364$.
\end{fullsolution}
\end{oefening}

\begin{oefening}[1]\label{oef:h3-rrsom}
\thema{rekenkundig: som}Bereken de som van de rekenkundige rij.
\begin{enumerate}
  \item $s_{15}$ als $u_1 = 8$ en $v = 3$
  \item $s_{11}$ als $u_1 = 40$ en $v = -6$
  \item $s_9$ als $u_1 = \tfrac12$ en $v = \tfrac14$
\end{enumerate}
\begin{solution}
\textbf{a.}~435 \quad \textbf{b.}~110 \quad \textbf{c.}~$\tfrac{27}{2}$
\end{solution}
\begin{fullsolution}
\textbf{a.}~$u_{15}=8+14\cdot 3=50$, dus $15\cdot\dfrac{8+50}{2}=435$.\\
\textbf{b.}~$u_{11}=40-60=-20$, dus $11\cdot\dfrac{40-20}{2}=110$.\\
\textbf{c.}~$u_9=\tfrac12+8\cdot\tfrac14=\tfrac52$, dus $9\cdot\dfrac{\tfrac12+\tfrac52}{2}=\tfrac{27}{2}$.
\end{fullsolution}
\end{oefening}

\begin{oefening}[1]\label{oef:h3-sigma}
\thema{somteken: uitschrijven}Schrijf uit en bereken.
\begin{reeks}[4]
  \item $\displaystyle\sum_{i=1}^{4} (2i + 1)$
  \item $\displaystyle\sum_{k=2}^{5} k^2$
  \item $\displaystyle\sum_{j=0}^{3} 2^{\,j}$
  \item $\displaystyle\sum_{i=1}^{5} (-1)^i$
\end{reeks}
\begin{solution}
\textbf{a.}~24 \quad \textbf{b.}~54 \quad \textbf{c.}~15 \quad \textbf{d.}~$-1$
\end{solution}
\begin{fullsolution}
\textbf{a.}~$3+5+7+9=24$.\\
\textbf{b.}~$4+9+16+25=54$.\\
\textbf{c.}~$1+2+4+8=15$.\\
\textbf{d.}~$-1+1-1+1-1=-1$.
\end{fullsolution}
\end{oefening}

\begin{oefening}[1]\label{oef:h3-mrsom}
\thema{meetkundig: som}Bereken de som van de meetkundige rij.
\begin{enumerate}
  \item $s_8$ als $u_1 = 5$ en $q = 2$
  \item $s_5$ als $u_1 = 81$ en $q = \tfrac13$
  \item $s_5$ als $u_1 = -2$ en $q = -3$
  \item $s_{20}$ als $u_1 = 7$ en $q = 1$
\end{enumerate}
\begin{solution}
\textbf{a.}~1275 \quad \textbf{b.}~121 \quad \textbf{c.}~$-122$ \quad \textbf{d.}~140
\end{solution}
\begin{fullsolution}
\textbf{a.}~$5\cdot(2^8-1)=1275$.\\
\textbf{b.}~$81\cdot\dfrac{(\tfrac13)^5-1}{\tfrac13-1}=81\cdot\dfrac{-\tfrac{242}{243}}{-\tfrac23}=121$.\\
\textbf{c.}~$-2\cdot\dfrac{(-3)^5-1}{-3-1}=-2\cdot\dfrac{-244}{-4}=-122$.\\
\textbf{d.}~$q=1$: alle termen zijn 7, dus $20\cdot 7=140$.
\end{fullsolution}
\end{oefening}

\begin{oefening}[1]\label{oef:h3-aantal}
\thema{aantal termen}Bereken de som.
\begin{enumerate}
  \item $11 + 14 + 17 + \cdots + 95$
  \item $6 + 12 + 24 + \cdots + 1536$
  \item $100 + 96 + 92 + \cdots + 4$
\end{enumerate}
\begin{solution}
\textbf{a.}~1537 \quad \textbf{b.}~3066 \quad \textbf{c.}~1300
\end{solution}
\begin{fullsolution}
\textbf{a.}~$v=3$: $11+(n-1)\cdot 3=95$ geeft $n=29$. Som: $29\cdot\dfrac{11+95}{2}=1537$.\\
\textbf{b.}~$q=2$: $6\cdot 2^{n-1}=1536$ geeft $2^{n-1}=256$, dus $n=9$. Som: $6\cdot(2^9-1)=3066$.\\
\textbf{c.}~$v=-4$: $100+(n-1)\cdot(-4)=4$ geeft $n=25$. Som: $25\cdot\dfrac{100+4}{2}=1300$.
\end{fullsolution}
\end{oefening}

\begin{oefening}[1]
\thema{fouten: exponent}Iemand schrijft: `$u_1 = 3$ en $q = 2$, dus $s_5 = 3 \cdot \dfrac{2^4 - 1}{2 - 1} = 45$.' Wat is er fout? Geef het juiste antwoord.
\begin{solution}
De exponent is $n$, niet $n-1$: $s_5=93$.
\end{solution}
\begin{fullsolution}
In de somformule staat de exponent $n$: $s_5=3\cdot(2^5-1)=93$. Controle: $3+6+12+24+48=93$.
\end{fullsolution}
\end{oefening}

\begin{oefening}[1]\label{oef:h3-noteer}
\thema{somteken: noteren}Noteer met het somteken.
\begin{enumerate}
  \item $5 + 10 + 15 + \cdots + 200$
  \item $1 + 4 + 9 + \cdots + 144$
  \item $2 + 4 + 8 + \cdots + 1024$
  \item $1 - 2 + 3 - 4 + \cdots - 20$
\end{enumerate}
\begin{solution}
\textbf{a.}~$\displaystyle\sum_{i=1}^{40} 5i$ \quad \textbf{b.}~$\displaystyle\sum_{i=1}^{12} i^2$ \quad \textbf{c.}~$\displaystyle\sum_{i=1}^{10} 2^{\,i}$ \quad \textbf{d.}~$\displaystyle\sum_{i=1}^{20} (-1)^{i+1} i$
\end{solution}
\begin{fullsolution}
\textbf{a.}~De veelvouden van 5, en $200=5\cdot 40$: $\displaystyle\sum_{i=1}^{40} 5i$.\\
\textbf{b.}~De kwadraten, en $144=12^2$: $\displaystyle\sum_{i=1}^{12} i^2$.\\
\textbf{c.}~De machten van 2, en $1024=2^{10}$: $\displaystyle\sum_{i=1}^{10} 2^{\,i}$.\\
\textbf{d.}~Positief bij oneven $i$: $\displaystyle\sum_{i=1}^{20} (-1)^{i+1} i$.\\
Andere schrijfwijzen kunnen ook, bijvoorbeeld met een andere beginwaarde.
\end{fullsolution}
\end{oefening}

\begin{oefening}[1]
\thema{fouten: aantal termen}Iemand schrijft: `$12 + 15 + \cdots + 60$ telt $\dfrac{60 - 12}{3} = 16$ termen.' Wat is er fout? Bereken de som.
\begin{solution}
Dat zijn 16 sprongen, dus 17 termen; de som is 612.
\end{solution}
\begin{fullsolution}
$\dfrac{60-12}{3}=16$ is het aantal sprongen. Er zijn dus 17 termen, en de som is $17\cdot\dfrac{12+60}{2}=612$.
\end{fullsolution}
\end{oefening}

\begin{oefening}[1]
\begin{metprent}{fiets}
\thema{context: rekenkundig}Een wielrenner fietst op de eerste dag van zijn trainingsperiode 20 km. Elke dag fietst hij 3 km meer dan de dag ervoor. Hoeveel kilometer fietst hij in 30 dagen?
\end{metprent}
\begin{solution}
1905 km
\end{solution}
\begin{fullsolution}
$u_{30}=20+29\cdot 3=107$, dus $s_{30}=30\cdot\dfrac{20+107}{2}=1905$.
\end{fullsolution}
\end{oefening}

\begin{oefening}[1]
\thema{de som als rij}Van een rij is $s_1 = 4$, $s_2 = 10$ en $s_3 = 18$. Bepaal $u_1$, $u_2$ en $u_3$.
\begin{solution}
4, 6, 8
\end{solution}
\begin{fullsolution}
$u_1=s_1=4$, $u_2=s_2-s_1=6$ en $u_3=s_3-s_2=8$.
\end{fullsolution}
\end{oefening}

\begin{oefening}[1]
\thema{somteken: berekenen}Bereken.
\begin{reeks}[2]
  \item $\displaystyle\sum_{i=1}^{30} (4i - 1)$
  \item $\displaystyle\sum_{i=1}^{6} 3 \cdot 2^{\,i}$
\end{reeks}
\begin{solution}
\textbf{a.}~1830 \quad \textbf{b.}~378
\end{solution}
\begin{fullsolution}
\textbf{a.}~$3+7+\cdots+119$: 30 termen, dus $30\cdot\dfrac{3+119}{2}=1830$.\\
\textbf{b.}~$6+12+\cdots+192$: $u_1=6$, $q=2$, dus $6\cdot(2^6-1)=378$.
\end{fullsolution}
\end{oefening}

\begin{oefening}[1]
\thema{rekenkundig: voorbeeld uit hoofdstuk 1}In hoofdstuk 1 telde trap $n$ precies $1 + 2 + \cdots + n$ blokjes. Toon met de somformule aan dat dit $\dfrac{n(n+1)}{2}$ is.
\begin{solution}
$u_1=1$, $u_n=n$, $n$ termen.
\end{solution}
\begin{fullsolution}
Dit is een rekenkundige rij met $u_1=1$, $u_n=n$ en $n$ termen, dus $s_n=n\cdot\dfrac{1+n}{2}=\dfrac{n(n+1)}{2}$.
\end{fullsolution}
\end{oefening}

\begin{oefening}[2]\label{oef:h3-somindex}
\thema{somteken: andere beginwaarde}Bereken.
\begin{reeks}
  \item $\displaystyle\sum_{k=5}^{24} (3k + 2)$
  \item $\displaystyle\sum_{j=3}^{8} 2^{\,j}$
  \item $\displaystyle\sum_{i=0}^{9} (10 - i)$
\end{reeks}
\begin{solution}
\textbf{a.}~910 \quad \textbf{b.}~504 \quad \textbf{c.}~55
\end{solution}
\begin{fullsolution}
\textbf{a.}~$17+20+\cdots+74$: $24-5+1=20$ termen, dus $20\cdot\dfrac{17+74}{2}=910$.\\
\textbf{b.}~$8+16+\cdots+256$: $8-3+1=6$ termen, dus $8\cdot(2^6-1)=504$.\\
\textbf{c.}~$10+9+\cdots+1$: $9-0+1=10$ termen, dus $10\cdot\dfrac{10+1}{2}=55$.
\end{fullsolution}
\end{oefening}

\begin{oefening}[2]\label{oef:h3-nuitsom}
\thema{aantal termen uit de som}Hoeveel termen van $2,\ 5,\ 8,\ \ldots$ moet je optellen om 155 te krijgen?
\begin{solution}
10
\end{solution}
\begin{fullsolution}
$u_n=3n-1$, dus $s_n=n\cdot\dfrac{2+3n-1}{2}=\dfrac{n(3n+1)}{2}$. Uit $3n^2+n-310=0$ volgt $n=10$ of $n=-\tfrac{31}{3}$. Enkel $n=10$ kan.
\end{fullsolution}
\end{oefening}

\begin{oefening}[2]\label{oef:h3-mrcontext}
\thema{context: meetkundig}Een filmpje krijgt op de eerste dag 200 weergaven. Elke dag zijn er 1,5 keer zoveel weergaven als de dag ervoor.
\begin{enumerate}
  \item Hoeveel weergaven zijn er in totaal na 10 dagen?
  \item Na hoeveel dagen is het totaal voor het eerst meer dan 100\,000?
\end{enumerate}
\begin{solution}
\textbf{a.}~ongeveer 22\,666 \quad \textbf{b.}~na 14 dagen
\end{solution}
\begin{fullsolution}
\textbf{a.}~$s_{10}=200\cdot\dfrac{1{,}5^{10}-1}{0{,}5}\approx 22\,666$.\\
\textbf{b.}~$400\,(1{,}5^{\,n}-1)>100\,000$ geeft $1{,}5^{\,n}>251$ en $n>13{,}63$, dus na 14 dagen.
\end{fullsolution}
\end{oefening}

\begin{oefening}[2]\label{oef:h3-uitsom}
\thema{de som als rij}Bepaal $u_n$. Is de rij rekenkundig, meetkundig of geen van beide?
\begin{enumerate}
  \item $s_n = 2n^2 + n$
  \item $s_n = 3^n - 1$
  \item $s_n = n^2 + 1$
\end{enumerate}
\begin{solution}
\textbf{a.}~$u_n=4n-1$, rekenkundig \quad \textbf{b.}~$u_n=2\cdot 3^{n-1}$, meetkundig \quad \textbf{c.}~$u_1=2$ en $u_n=2n-1$ voor $n\geqslant 2$: geen van beide
\end{solution}
\begin{fullsolution}
\textbf{a.}~$s_{n-1}=2n^2-3n+1$, dus $u_n=4n-1$. Ook $u_1=3$ past hierin.\\
\textbf{b.}~$u_n=3^n-3^{n-1}=3^{n-1}(3-1)=2\cdot 3^{n-1}$. Ook $u_1=2$ past hierin.\\
\textbf{c.}~$u_n=n^2+1-\bigl((n-1)^2+1\bigr)=2n-1$ voor $n\geqslant 2$. Maar $u_1=s_1=2$ past daar niet in: de rij is $2,\ 3,\ 5,\ 7,\ \ldots$
\end{fullsolution}
\end{oefening}

\begin{oefening}[2]\label{oef:h3-sparen}
\thema{sparen}Je stort 10 jaar lang aan het einde van elk jaar \eur 600 op een rekening met 2,5\% intrest.
\begin{enumerate}
  \item Hoeveel staat er na de laatste storting op?
  \item Hoeveel daarvan is intrest?
  \item Hoeveel had er gestaan als je de \eur 6000 meteen aan het begin had gestort?
\end{enumerate}
\begin{solution}
\textbf{a.}~ongeveer \eur 6722,03 \quad \textbf{b.}~ongeveer \eur 722,03 \quad \textbf{c.}~ongeveer \eur 7680,51
\end{solution}
\begin{fullsolution}
\textbf{a.}~$E=600\cdot\dfrac{1{,}025^{10}-1}{0{,}025}\approx 6722{,}03$.\\
\textbf{b.}~$6722{,}03-6000=722{,}03$.\\
\textbf{c.}~$6000\cdot 1{,}025^{10}\approx 7680{,}51$: het geld staat dan langer op de rekening.
\end{fullsolution}
\end{oefening}

\begin{oefening}[2]
\thema{rekenkundig: onbekende uit de som}Bepaal de onbekende van de rekenkundige rij.
\begin{enumerate}
  \item $v = 3$ en $s_{12} = 342$. Bepaal $u_1$.
  \item $u_1 = 5$ en $s_{10} = 185$. Bepaal $v$.
  \item $s_5 = 40$ en $s_{10} = 130$. Bepaal $u_1$ en $v$.
\end{enumerate}
\begin{solution}
\textbf{a.}~$u_1=12$ \quad \textbf{b.}~$v=3$ \quad \textbf{c.}~$u_1=4$, $v=2$
\end{solution}
\begin{fullsolution}
Gebruik $s_n=\tfrac{n}{2}\bigl(2u_1+(n-1)v\bigr)$.\\
\textbf{a.}~$6(2u_1+33)=342$ geeft $2u_1+33=57$, dus $u_1=12$.\\
\textbf{b.}~$5(10+9v)=185$ geeft $10+9v=37$, dus $v=3$.\\
\textbf{c.}~$\tfrac52(2u_1+4v)=40$ geeft $u_1+2v=8$. $5(2u_1+9v)=130$ geeft $2u_1+9v=26$. Dus $5v=10$, $v=2$ en $u_1=4$.
\end{fullsolution}
\end{oefening}

\begin{oefening}[2]
\begin{metprent}{klok}
\thema{context: rekenkundig}Een klok slaat elk heel uur het aantal uren, van 1 tot en met 12, en elk half uur één keer. Hoeveel slagen hoor je in één etmaal?
\end{metprent}
\begin{solution}
180
\end{solution}
\begin{fullsolution}
Elke twaalf uur: $1+2+\cdots+12=12\cdot\dfrac{1+12}{2}=78$ slagen. In een etmaal is dat $2\cdot 78=156$. Daar komen 24 halve uren bij: $156+24=180$.
\end{fullsolution}
\end{oefening}

\begin{oefening}[2]
\thema{meetkundig: onbekende uit de som}Bepaal de onbekende van de meetkundige rij.
\begin{enumerate}
  \item $q = 2$ en $s_6 = 189$. Bepaal $u_1$.
  \item $u_1 = 2$ en $s_3 = 14$. Bepaal $q$. Let op: er zijn twee oplossingen.
\end{enumerate}
\begin{solution}
\textbf{a.}~$u_1=3$ \quad \textbf{b.}~$q=2$ of $q=-3$
\end{solution}
\begin{fullsolution}
\textbf{a.}~$u_1\cdot(2^6-1)=63\,u_1=189$, dus $u_1=3$.\\
\textbf{b.}~$s_3=u_1+u_1q+u_1q^2$, dus $2+2q+2q^2=14$ en $q^2+q-6=0$. Dat geeft $q=2$ of $q=-3$. Controle: $2+4+8=14$ en $2-6+18=14$.
\end{fullsolution}
\end{oefening}

\begin{oefening}[2]
\thema{aantal termen: voorwaarde}Bereken de som van alle getallen van drie cijfers die bij deling door 7 als rest 2 geven.
\begin{solution}
70\,692
\end{solution}
\begin{fullsolution}
Het kleinste is $100=7\cdot 14+2$ en het grootste is $996=7\cdot 142+2$. Dat zijn $142-14+1=129$ termen, dus de som is $129\cdot\dfrac{100+996}{2}=70\,692$.
\end{fullsolution}
\end{oefening}

\begin{oefening}[2]
\thema{fouten: geen rekenkundige rij}Iemand berekent $1 + 4 + 9 + 16 + \cdots + 100$ als $10 \cdot \dfrac{1 + 100}{2} = 505$. Wat is er fout? Wat is de juiste uitkomst?
\begin{solution}
De rij is niet rekenkundig; de som is 385.
\end{solution}
\begin{fullsolution}
De verschillen 3, 5, 7 zijn niet gelijk, dus de rij is niet rekenkundig en de formule geldt niet. Optellen geeft $1+4+9+16+25+36+49+64+81+100=385$.
\end{fullsolution}
\end{oefening}

\begin{oefening}[2]
\thema{meetkundig: wisselend teken}Bereken $1 - \tfrac12 + \tfrac14 - \tfrac18 + \cdots - \tfrac{1}{128}$.
\begin{solution}
$\tfrac{85}{128}$
\end{solution}
\begin{fullsolution}
$q=-\tfrac12$ en $\tfrac{1}{128}=(\tfrac12)^7$, dus $n-1=7$ en $n=8$. De som is $\dfrac{(-\tfrac12)^8-1}{-\tfrac12-1}=\dfrac{-\tfrac{255}{256}}{-\tfrac32}=\dfrac{85}{128}$.
\end{fullsolution}
\end{oefening}

\begin{oefening}[2]
\thema{somteken: splitsen}Bereken $\displaystyle\sum_{i=1}^{10} \left(2i + 2^{\,i}\right)$.
\begin{solution}
2156
\end{solution}
\begin{fullsolution}
Splits de som: $\displaystyle\sum_{i=1}^{10} 2i=10\cdot\dfrac{2+20}{2}=110$ en $\displaystyle\sum_{i=1}^{10} 2^{\,i}=2\cdot(2^{10}-1)=2046$. Samen 2156.
\end{fullsolution}
\end{oefening}

\begin{oefening}[2]
\thema{context: meetkundig}Een bal valt van 3 m hoogte en stuit telkens terug tot 80\% van de vorige hoogte. Welke afstand heeft de bal afgelegd op het moment dat hij voor de zesde keer de grond raakt?
\begin{solution}
Ongeveer 19,14 m
\end{solution}
\begin{fullsolution}
De bal valt 3 m, en daarna gaat hij 5 keer omhoog en weer omlaag. De stuithoogtes zijn $2{,}4;\ 1{,}92;\ \ldots$ met $q=0{,}8$. Hun som is $2{,}4\cdot\dfrac{0{,}8^5-1}{0{,}8-1}=8{,}06784$. De afstand is $3+2\cdot 8{,}06784\approx 19{,}14$ m.
\end{fullsolution}
\end{oefening}

\begin{oefening}[2]
\thema{meetkundig: aantal termen uit de som}Hoeveel termen van $3,\ 6,\ 12,\ \ldots$ moet je minstens optellen om meer dan een miljoen te krijgen?
\begin{solution}
19
\end{solution}
\begin{fullsolution}
$3\,(2^{\,n}-1)>1\,000\,000$ geeft $2^{\,n}>333\,334{,}33\ldots$ en $n>18{,}35$, dus $n=19$.
\end{fullsolution}
\end{oefening}

\begin{oefening}[2]
\thema{rekenkundig: grootste som}Je telt de termen van $61,\ 57,\ 53,\ \ldots$ één voor één op. Bij hoeveel termen is de som het grootst? Hoe groot is ze dan?
\begin{solution}
Bij 16 termen; 496
\end{solution}
\begin{fullsolution}
De som groeit zolang de termen positief zijn. $u_n=65-4n>0$ geeft $n<16{,}25$, dus $u_{16}=1$ is de laatste positieve term. Dan is $s_{16}=16\cdot\dfrac{61+1}{2}=496$.
\end{fullsolution}
\end{oefening}

\begin{oefening}[2]
\thema{context: meetkundig}Je mag kiezen: \eur 100\,000 ineens, of 30 dagen lang geld volgens dit schema: 1 cent op dag 1, 2 cent op dag 2, 4 cent op dag 3, enzovoort. Wat kies je?
\begin{solution}
Het schema: \eur 10\,737\,418,23
\end{solution}
\begin{fullsolution}
$s_{30}=1\cdot(2^{30}-1)=1\,073\,741\,823$ cent, dat is \eur 10\,737\,418,23. Dat is ruim 100 keer meer. Na 20 dagen heb je nochtans maar $2^{20}-1$ cent, ongeveer \eur 10\,486.
\end{fullsolution}
\end{oefening}

\begin{oefening}[3]
\thema{rekenkundig: bewijs}Bewijs: als $s_n = an^2 + bn$ voor alle $n$, met $a$ en $b$ vaste getallen, dan is de rij rekenkundig. Wat is het verschil?
\begin{solution}
Het verschil is $2a$.
\end{solution}
\begin{fullsolution}
$s_{n-1}=a(n-1)^2+b(n-1)=an^2-2an+a+bn-b$, dus $u_n=s_n-s_{n-1}=2an-a+b$ voor $n\geqslant 2$. Ook $u_1=s_1=a+b$ past hierin. De coëfficiënt van $n$ is $2a$, dus $v=2a$.
\end{fullsolution}
\end{oefening}

\begin{oefening}[3]
\begin{metprent}{slak}
\thema{meetkundig: de som loopt niet op}Een slak kruipt op dag 1 een meter. Elke dag kruipt ze de helft van de dag ervoor.
\begin{enumerate}
  \item Toon aan dat de afstand na $n$ dagen $s_n = 2 - 2 \cdot (\tfrac12)^n$ meter is.
  \item Na hoeveel dagen heeft de slak voor het eerst meer dan 1,99 m afgelegd?
  \item Haalt de slak ooit 2 m? Leg uit.
\end{enumerate}
\end{metprent}
\begin{solution}
\textbf{b.}~na 8 dagen \quad \textbf{c.}~nee
\end{solution}
\begin{fullsolution}
\textbf{a.}~$s_n=\dfrac{(\tfrac12)^n-1}{\tfrac12-1}=-2\bigl((\tfrac12)^n-1\bigr)=2-2\cdot(\tfrac12)^n$.\\
\textbf{b.}~$2\cdot(\tfrac12)^n<0{,}01$ geeft $(\tfrac12)^n<0{,}005$ en $n>7{,}64$, dus na 8 dagen.\\
\textbf{c.}~Nee: $2\cdot(\tfrac12)^n$ is altijd positief, dus $s_n<2$. De som komt wel zo dicht bij 2 als je wilt. In hoofdstuk 4 noemen we dat een limiet.
\end{fullsolution}
\end{oefening}

\begin{oefening}[3]
\thema{rekenkundig: slim groeperen}Bereken $100^2 - 99^2 + 98^2 - 97^2 + \cdots + 2^2 - 1^2$ zonder rekentoestel.
\begin{solution}
5050
\end{solution}
\begin{fullsolution}
Gebruik $a^2-b^2=(a-b)(a+b)$: elk paar is $(a+b)$ met $a-b=1$. Dat geeft $199+195+\cdots+3$. Dit is een rekenkundige rij met 50 termen, dus de som is $50\cdot\dfrac{199+3}{2}=5050$.
\end{fullsolution}
\end{oefening}

\begin{oefening}[3]
\thema{lening}Je leent \eur 10\,000 aan 4\% per jaar, en betaalt 5 jaar lang aan het einde van elk jaar een vast bedrag $a$. Na de laatste betaling is de lening volledig terugbetaald.
\begin{enumerate}
  \item Waarom geldt $a \cdot \dfrac{1{,}04^5 - 1}{1{,}04 - 1} = 10\,000 \cdot 1{,}04^5$? Denk aan twee rekeningen die na 5 jaar evenveel waard moeten zijn.
  \item Bereken $a$. Hoeveel intrest betaal je in totaal?
\end{enumerate}
\begin{solution}
\textbf{b.}~$a\approx$ \eur 2246,27; ongeveer \eur 1231,36 intrest
\end{solution}
\begin{fullsolution}
\textbf{a.}~Links: de betalingen, elk aangegroeid tot het einde van jaar 5. Rechts: de schuld, die zonder betalingen tot dan aangroeit. Die twee moeten gelijk zijn.\\
\textbf{b.}~$a=10\,000\cdot 1{,}04^5\cdot\dfrac{0{,}04}{1{,}04^5-1}\approx 2246{,}27$. In totaal betaal je $5a\approx 11\,231{,}36$, dus ongeveer \eur 1231,36 intrest.
\end{fullsolution}
\end{oefening}

\begin{oefening}[3]
\thema{som in een tabel}Je telt alle getallen op in de tafels van vermenigvuldiging van 1 tot en met 10, dus de $10 \times 10$ producten van $1 \cdot 1$ tot en met $10 \cdot 10$.
\begin{enumerate}
  \item Hoeveel is de som van de tafel van $k$?
  \item Toon aan dat het totaal gelijk is aan $(1 + 2 + \cdots + 10)^2$ en bereken het.
\end{enumerate}
\begin{solution}
\textbf{a.}~$55k$ \quad \textbf{b.}~3025
\end{solution}
\begin{fullsolution}
\textbf{a.}~$k\cdot 1+k\cdot 2+\cdots+k\cdot 10=k\,(1+2+\cdots+10)=55k$.\\
\textbf{b.}~Het totaal is $55\cdot 1+55\cdot 2+\cdots+55\cdot 10=55\cdot(1+2+\cdots+10)=55^2=3025$.
\end{fullsolution}
\end{oefening}

\begin{oefening}[3]
\thema{bewijs: rekenkundige som}Bewijs: bij een rekenkundige rij is $s_n = \dfrac{n}{2}\bigl(2u_1 + (n-1)v\bigr)$. Toon daarmee aan dat de som van de eerste $n$ oneven getallen gelijk is aan $n^2$.
\begin{solution}
$u_1=1$, $v=2$ invullen geeft $n^2$.
\end{solution}
\begin{fullsolution}
Vul $u_n=u_1+(n-1)v$ in de somformule in: $s_n=n\cdot\dfrac{u_1+u_1+(n-1)v}{2}=\dfrac{n}{2}\bigl(2u_1+(n-1)v\bigr)$.\\
Voor de oneven getallen is $u_1=1$ en $v=2$, dus $s_n=\dfrac{n}{2}\bigl(2+2(n-1)\bigr)=\dfrac{n}{2}\cdot 2n=n^2$.
\end{fullsolution}
\end{oefening}
\end{oefeningen}

%----------------------------------------------------------------------------------------
% KOPPEL: verbanden met andere leerstof
%----------------------------------------------------------------------------------------
\opdrachtsectie{Koppel}

\lkr{Matrices en complexe getallen zijn gezien in het 5de jaar. Logica, kansrekenen en integralen staan ook in het 6de jaar: kijk na of ze al aan bod kwamen, of laat die oefeningen weg.}

\begin{oefeningen}
\begin{oefening}[3]\label{oef:h3-inductie}
\koppel{discreet}{Logica en bewijzen} Voor kwadraten bestaat geen rekenkundige of meetkundige som, maar wel een formule: $\displaystyle\sum_{i=1}^{n} i^2 = \dfrac{n(n+1)(2n+1)}{6}$.
\begin{enumerate}
  \item Controleer de formule voor $n = 10$.
  \item Bewijs de formule met volledige inductie.
\end{enumerate}
\begin{solution}
\textbf{a.}~385
\end{solution}
\begin{fullsolution}
\textbf{a.}~$\dfrac{10\cdot 11\cdot 21}{6}=385$.\\
\textbf{b.}~Basis: voor $n=1$ staat er $1=\dfrac{1\cdot 2\cdot 3}{6}$.\\
Stap: stel dat de formule geldt voor $n=k$. Dan is de som tot $k+1$ gelijk aan\\
$\dfrac{k(k+1)(2k+1)}{6}+(k+1)^2=\dfrac{(k+1)\bigl(2k^2+k+6k+6\bigr)}{6}=\dfrac{(k+1)(k+2)(2k+3)}{6}$.\\
Dat is de formule voor $n=k+1$.
\end{fullsolution}
\end{oefening}

\begin{oefening}[2]
\koppel{data}{Kansrekenen} Je gooit met een dobbelsteen tot je een zes gooit. $u_i$ is de kans dat de eerste zes valt bij worp $i$.
\begin{enumerate}
  \item Bereken met de somformule de kans dat de eerste zes valt in de eerste $n$ worpen.
  \item Bereken diezelfde kans via de complementaire gebeurtenis. Komt hetzelfde uit?
  \item Hoeveel worpen heb je nodig om minstens 90\% kans op een zes te hebben?
\end{enumerate}
\begin{solution}
\textbf{a.}~$1-(\tfrac56)^n$ \quad \textbf{b.}~ja \quad \textbf{c.}~13
\end{solution}
\begin{fullsolution}
\textbf{a.}~$u_i=\tfrac16\cdot(\tfrac56)^{i-1}$, dus $s_n=\tfrac16\cdot\dfrac{(\tfrac56)^n-1}{\tfrac56-1}=1-(\tfrac56)^n$.\\
\textbf{b.}~De kans op $n$ keer geen zes is $(\tfrac56)^n$, dus de gevraagde kans is $1-(\tfrac56)^n$.\\
\textbf{c.}~$(\tfrac56)^n\leqslant 0{,}1$ geeft $n\geqslant 12{,}63$, dus 13 worpen.
\end{fullsolution}
\end{oefening}

\begin{oefening}[2]
\koppel{algebra}{Complexe getallen} De somformule geldt ook als $q$ een complex getal is.
\begin{enumerate}
  \item Bereken $1 + \mathrm{i} + \mathrm{i}^2 + \cdots + \mathrm{i}^{10}$ met de somformule, en controleer door de termen uit te schrijven.
  \item Stel dat $w \neq 1$ en $w^n = 1$. Toon aan dat $1 + w + w^2 + \cdots + w^{n-1} = 0$. Wat betekent dat voor de $n$-de wortels uit 1?
\end{enumerate}
\begin{solution}
\textbf{a.}~$\mathrm{i}$ \quad \textbf{b.}~Hun som is 0.
\end{solution}
\begin{fullsolution}
\textbf{a.}~$\dfrac{\mathrm{i}^{11}-1}{\mathrm{i}-1}=\dfrac{-\mathrm{i}-1}{\mathrm{i}-1}=\dfrac{(-\mathrm{i}-1)(-\mathrm{i}-1)}{(\mathrm{i}-1)(-\mathrm{i}-1)}=\dfrac{2\mathrm{i}}{2}=\mathrm{i}$. Controle: per vier termen is de som 0, en $1+\mathrm{i}-1=\mathrm{i}$ blijft over.\\
\textbf{b.}~$\dfrac{w^n-1}{w-1}=\dfrac{0}{w-1}=0$. Neem voor $w$ de wortel met het kleinste positieve argument: dan zijn $1,\ w,\ \ldots,\ w^{n-1}$ precies alle $n$-de wortels uit 1. Hun som is dus 0 (voor $n\geqslant 2$).
\end{fullsolution}
\end{oefening}

\begin{oefening}[2]
\koppel{algebra}{Matrices} Gegeven: $A = \begin{pmatrix} 2 & 0 \\ 0 & 3 \end{pmatrix}$ en $I = \begin{pmatrix} 1 & 0 \\ 0 & 1 \end{pmatrix}$.
\begin{enumerate}
  \item Bereken $S = I + A + A^2 + A^3$. Welke twee sommen herken je?
  \item Toon aan: $(A - I)\,S = A^4 - I$. Vergelijk met het bewijs van de somformule.
\end{enumerate}
\begin{solution}
\textbf{a.}~$S=\begin{pmatrix} 15 & 0 \\ 0 & 40 \end{pmatrix}$
\end{solution}
\begin{fullsolution}
\textbf{a.}~$S=\begin{pmatrix} 15 & 0 \\ 0 & 40 \end{pmatrix}$: op de diagonaal staan $1+2+4+8$ en $1+3+9+27$.\\
\textbf{b.}~$(A-I)S=AS-S=A+A^2+A^3+A^4-I-A-A^2-A^3=A^4-I$. Dit is dezelfde redenering als $q\cdot s_n-s_n$. Controle: $\begin{pmatrix} 1 & 0 \\ 0 & 2 \end{pmatrix}\begin{pmatrix} 15 & 0 \\ 0 & 40 \end{pmatrix}=\begin{pmatrix} 15 & 0 \\ 0 & 80 \end{pmatrix}$.
\end{fullsolution}
\end{oefening}

\begin{oefening}[3]
\koppel{analyse}{Integralen} Verdeel $[0, 1]$ in $n$ gelijke stukken en teken onder de grafiek van $f(x) = x^2$ de rechthoeken met breedte $\tfrac1n$ en als hoogte de functiewaarde rechts in elk stuk.

\begin{tikzpicture}[x=2.6cm,y=1.7cm,font=\scriptsize]
  \foreach \i in {1,2,3,4,5}{
    \draw[accent,fill=accent!15] ({(\i-1)/5},0) rectangle ({\i/5},{(\i/5)^2});}
  \draw[->,lijngrijs] (-.05,0) -- (1.15,0) node[right,text=black] {$x$};
  \draw[->,lijngrijs] (0,-.05) -- (0,1.15) node[above,text=black] {$y$};
  \draw[thick,domain=0:1.04,samples=30] plot (\x,{\x*\x});
  \node[below=1pt,text=grijs] at (1,0) {1};
\end{tikzpicture}
\begin{enumerate}
  \item Toon aan dat de totale oppervlakte van de rechthoeken gelijk is aan $\dfrac{1}{n^3} \displaystyle\sum_{i=1}^{n} i^2$.
  \item Gebruik de formule uit oefening \ref{oef:h3-inductie} en bereken die oppervlakte voor $n = 10$ en $n = 100$. Welk getal benader je?
\end{enumerate}
\begin{solution}
\textbf{b.}~0,385 en 0,33835; de benadering gaat naar $\tfrac13$
\end{solution}
\begin{fullsolution}
\textbf{a.}~De $i$-de rechthoek heeft breedte $\tfrac1n$ en hoogte $(\tfrac{i}{n})^2$, dus oppervlakte $\dfrac{i^2}{n^3}$.\\
\textbf{b.}~$n=10$: $\dfrac{385}{1000}=0{,}385$. $n=100$: $\dfrac{338\,350}{1\,000\,000}=0{,}33835$. Het getal is $\displaystyle\int_0^1 x^2\,\mathrm{d}x=\tfrac13$.
\end{fullsolution}
\end{oefening}
\end{oefeningen}

%----------------------------------------------------------------------------------------
\opdrachtsectie{Test jezelf}

\begin{oefeningen}
\begin{oefening}[1]\label{oef:h3-t1}
Een rekenkundige rij heeft $u_1 = 6$ en $v = 5$. Bereken $s_{12}$.
\begin{solution}
402
\end{solution}
\begin{fullsolution}
$u_{12}=6+55=61$, dus $s_{12}=12\cdot\dfrac{6+61}{2}=402$.
\end{fullsolution}
\end{oefening}

\begin{oefening}[1]\label{oef:h3-t2}
Bereken $13 + 17 + 21 + \cdots + 93$.
\begin{solution}
1113
\end{solution}
\begin{fullsolution}
$13+(n-1)\cdot 4=93$ geeft $n=21$, dus de som is $21\cdot\dfrac{13+93}{2}=1113$.
\end{fullsolution}
\end{oefening}

\begin{oefening}[1]\label{oef:h3-t3}
Bereken $s_7$ van de rij $4,\ 12,\ 36,\ \ldots$
\begin{solution}
4372
\end{solution}
\begin{fullsolution}
$q=3$, dus $s_7=4\cdot\dfrac{3^7-1}{3-1}=4372$.
\end{fullsolution}
\end{oefening}

\begin{oefening}[2]\label{oef:h3-t4}
Bereken $\displaystyle\sum_{k=4}^{13} (3k - 2)$.
\begin{solution}
235
\end{solution}
\begin{fullsolution}
$10+13+\cdots+37$: $13-4+1=10$ termen, dus $10\cdot\dfrac{10+37}{2}=235$.
\end{fullsolution}
\end{oefening}

\begin{oefening}[1]\label{oef:h3-t5}
Noteer met het somteken: $7 + 14 + 21 + \cdots + 84$.
\begin{solution}
$\displaystyle\sum_{i=1}^{12} 7i$
\end{solution}
\begin{fullsolution}
De veelvouden van 7, en $84=7\cdot 12$: $\displaystyle\sum_{i=1}^{12} 7i$.
\end{fullsolution}
\end{oefening}

\begin{oefening}[2]\label{oef:h3-t6}
Hoeveel termen van $5,\ 9,\ 13,\ \ldots$ moet je optellen om 230 te krijgen?
\begin{solution}
10
\end{solution}
\begin{fullsolution}
$u_n=4n+1$, dus $s_n=n\cdot\dfrac{5+4n+1}{2}=n(2n+3)$. Uit $2n^2+3n-230=0$ volgt $n=10$ of $n=-\tfrac{23}{2}$. Enkel $n=10$ kan.
\end{fullsolution}
\end{oefening}

\begin{oefening}[2]\label{oef:h3-t7}
Gegeven: $s_n = 3n^2 - n$. Bepaal $u_n$.
\begin{solution}
$u_n=6n-4$
\end{solution}
\begin{fullsolution}
$s_{n-1}=3(n-1)^2-(n-1)=3n^2-7n+4$, dus $u_n=6n-4$ voor $n\geqslant 2$. Ook $u_1=s_1=2$ past hierin.
\end{fullsolution}
\end{oefening}

\begin{oefening}[2]\label{oef:h3-t8}
Je loopt op dag 1 een afstand van 2 km, en elke dag 10\% meer dan de dag ervoor. Hoeveel kilometer loop je in 7 dagen?
\begin{solution}
Ongeveer 18,97 km
\end{solution}
\begin{fullsolution}
$s_7=2\cdot\dfrac{1{,}1^7-1}{0{,}1}\approx 18{,}97$.
\end{fullsolution}
\end{oefening}

\begin{oefening}[2]\label{oef:h3-t9}
Je stort 5 jaar lang aan het einde van elk jaar \eur 800 op een rekening met 3\% intrest. Hoeveel staat er na de laatste storting op?
\begin{solution}
Ongeveer \eur 4247,31
\end{solution}
\begin{fullsolution}
$E=800\cdot\dfrac{1{,}03^5-1}{0{,}03}\approx 4247{,}31$.
\end{fullsolution}
\end{oefening}
\end{oefeningen}

\herhaalschema{%
  \fv{oef:h3-t1}{vb:h3-rr}{oef:h3-rrsom}%
  \fv{oef:h3-t2}{vb:h3-aantal}{oef:h3-aantal}%
  \fv{oef:h3-t3}{vb:h3-mr}{oef:h3-mrsom}%
  \fv{oef:h3-t4}{vb:h3-somteken}{oef:h3-somindex}%
  \fv{oef:h3-t5}{vb:h3-somteken}{oef:h3-noteer}%
  \fv{oef:h3-t6}{vb:h3-nuit}{oef:h3-nuitsom}%
  \fv{oef:h3-t7}{vb:h3-uitsom}{oef:h3-uitsom}%
  \fv{oef:h3-t8}{vb:h3-mr}{oef:h3-mrcontext}%
  \fv{oef:h3-t9}{vb:h3-sparen}{oef:h3-sparen}}

\ifzelfstudie
\section*{Maak je eigen samenvatting}
Sluit dit boek en schrijf in je schrift: de somformule van een rekenkundige en van een meetkundige rij, met het bewijs van elk, hoe je het aantal termen vindt, wat het somteken betekent, hoe je $u_n$ uit $s_n$ haalt, en welke fouten je wilt vermijden. Vergelijk daarna met de kaders.
\fi
